% !TEX root =./ECE343_Practice_main.tex
%
% L35 practice set (Maxwell's Equations in a Lossless, Source-Free Region I):
% vector phasors, the phasor Ampere–Maxwell law, decoupling the magnetic field,
% the Helmholtz equation and its phase constant, and the line analogy.

\fullwidth{\textbf{Learning Objective 9.1: } I understand the process of
decoupling and finding solutions for Maxwell's equations for dynamic fields.
\\[4pt]\normalfont\small Show your work. Write a \textbf{Documentation} line at the
end (\textbf{None} if you did not collaborate).}

\begin{questions}

%%%%%%%%%%%%%%%%%%%%%%%%%%%%%%%%%%%%%%%%%
\question \textbf{Vector phasors.} In a lossless, source-free region, the electric
field is
\eq{ \mathbf{E}(z,t) = \hat{\mathbf{x}}\ 5\cos\left(\omega t - \beta z + 30^\circ\right) + \hat{\mathbf{y}}\ 3\sin\left(\omega t - \beta z\right)\ \text{V/m} }
\begin{parts}

	\part Write the vector phasor $\widetilde{\mathbf{E}}(z)$.
	\begin{solution}[1.2in]
	Write the sine as a cosine, $\sin\theta = \cos\left(\theta - 90^\circ\right)$, so
	each component is a cosine with its own amplitude and phase:
	\eq{
		\widetilde{\mathbf{E}}(z) &= \hat{\mathbf{x}}\ 5e^{j30^\circ}e^{-j\beta z} + \hat{\mathbf{y}}\ 3e^{-j90^\circ}e^{-j\beta z} \\
		&= \left(\hat{\mathbf{x}}\ 5e^{j30^\circ} - j\ \hat{\mathbf{y}}\ 3\right)e^{-j\beta z}\ \text{V/m}
	}
	\end{solution}

	\part Starting from $\nabla\times\mathbf{H} = \mathbf{J} + \partial\mathbf{D}/\partial t$
	in the time domain, derive the phasor form of the Ampere–Maxwell law in this
	region. State each assumption where you use it.
	\begin{solution}[2.2in]
	Substitute $\mathbf{H} = \text{Re}\left[\widetilde{\mathbf{H}}e^{j\omega t}\right]$
	and $\mathbf{D} = \text{Re}\left[\widetilde{\mathbf{D}}e^{j\omega t}\right]$. The curl
	acts on position, so it passes through the real part and $e^{j\omega t}$; the time
	derivative acts only on $e^{j\omega t}$:
	\eq{
		\text{Re}\left[\left(\nabla\times\widetilde{\mathbf{H}}\right)e^{j\omega t}\right] &= \text{Re}\left[\left(\widetilde{\mathbf{J}} + j\omega\widetilde{\mathbf{D}}\right)e^{j\omega t}\right] \\
		\nabla\times\widetilde{\mathbf{H}} &= \widetilde{\mathbf{J}} + j\omega\widetilde{\mathbf{D}}
	}
	The equation holds at every instant, so the complex vectors are equal. The region is lossless,
	so $\sigma = 0$ and $\widetilde{\mathbf{J}} = \sigma\widetilde{\mathbf{E}} = 0$ (no
	impressed current either, since the region is source-free). A simple medium has
	$\widetilde{\mathbf{D}} = \epsilon\widetilde{\mathbf{E}}$:
	\eq{ \nabla\times\widetilde{\mathbf{H}} = j\omega\epsilon\widetilde{\mathbf{E}} }
	\end{solution}

\end{parts}

\newpage
%%%%%%%%%%%%%%%%%%%%%%%%%%%%%%%%%%%%%%%%%
\question \textbf{Decoupling the magnetic field.}
\begin{parts}

	\part Starting from the Ampere–Maxwell law, in phasor form, in a lossless,
	source-free region, derive the Helmholtz equation for the magnetic field. Name
	each of Maxwell's equations where you use it.
	\begin{solution}[2.4in]
	Take the curl of the Ampere–Maxwell law. The constant $j\omega\epsilon$ comes
	outside the curl:
	\eq{ \nabla\times\left(\nabla\times\widetilde{\mathbf{H}}\right) = j\omega\epsilon\ \nabla\times\widetilde{\mathbf{E}} }
	Faraday's law gives $\nabla\times\widetilde{\mathbf{E}} = -j\omega\mu\widetilde{\mathbf{H}}$.
	The identity $\nabla\times\nabla\times\mathbf{F} = \nabla\left(\nabla\cdot\mathbf{F}\right) - \nabla^2\mathbf{F}$
	expands the left side, and Gauss's law for magnetism gives
	$\nabla\cdot\widetilde{\mathbf{H}} = 0$:
	\eq{
		\nabla\left(\nabla\cdot\widetilde{\mathbf{H}}\right) - \nabla^2\widetilde{\mathbf{H}} &= j\omega\epsilon\left(-j\omega\mu\widetilde{\mathbf{H}}\right) = \omega^2\mu\epsilon\ \widetilde{\mathbf{H}} \\
		\nabla^2\widetilde{\mathbf{H}} + \beta^2\widetilde{\mathbf{H}} &= 0, \qquad \beta = \omega\sqrt{\mu\epsilon}
	}
	Gauss's law for the electric field is the one equation this derivation does not
	use; the derivation of the electric-field equation uses it instead.
	\end{solution}

	\part On a transmission line, which line quantities play the parts of $\mu$ and
	$\epsilon$, and which plays the part of $\widetilde{\mathbf{H}}$?
	\begin{solution}[0.9in]
	The inductance per meter $L^{\prime}$ plays $\mu$, the capacitance per meter
	$C^{\prime}$ plays $\epsilon$, and the current $\widetilde{I}$ plays
	$\widetilde{\mathbf{H}}$, as the voltage $\widetilde{V}$ plays $\widetilde{\mathbf{E}}$.
	\end{solution}

\end{parts}

\newpage
%%%%%%%%%%%%%%%%%%%%%%%%%%%%%%%%%%%%%%%%%
\question \textbf{The Helmholtz equation.} A lossless dielectric has $\epsilon_r = 4$
and $\mu_r = 1$. The frequency is $300\ \mhz$.
\begin{parts}

	\begin{minipage}[t]{0.58\linewidth}
		\part Find $\beta$, the wavelength, and the phase velocity.
		\begin{solution}[1.3in]
		\eq{
			\beta &= \frac{\omega\sqrt{\epsilon_r}}{c} = \frac{2\pi\times3\times10^{8}\times2}{2.998\times10^{8}} = 12.57\ \text{rad/m} \\
			\lambda &= \frac{2\pi}{\beta} = 0.500\ \m, \qquad v_p = \frac{c}{\sqrt{\epsilon_r}} = 1.50\times10^{8}\ \text{m/s}
		}
		\end{solution}
	\end{minipage}\hfill%
	\begin{minipage}[t]{0.37\linewidth}
		\raggedleft \vspace{12pt}
		$\beta$ = \ansbox[1.3in]{$12.57$ rad/m}

		\vspace{4pt}
		$\lambda$ = \ansbox[1.3in]{$0.500\ \m$}

		\vspace{4pt}
		$v_p$ = \ansbox[1.3in]{$1.50\times10^{8}$ m/s}
	\end{minipage}

	\part Show that $\widetilde{\mathbf{E}} = \hat{\mathbf{x}}\ E_0e^{-j\beta z}$
	satisfies both the Helmholtz equation and $\nabla\cdot\widetilde{\mathbf{E}} = 0$.
	\begin{solution}[1.2in]
	Only the $x$ component is nonzero, and it depends only on $z$:
	\eq{
		\nabla^2\widetilde{\mathbf{E}} &= \hat{\mathbf{x}}\ \frac{d^2}{dz^2}\left(E_0e^{-j\beta z}\right) = \hat{\mathbf{x}}\left(-j\beta\right)^2E_0e^{-j\beta z} = -\beta^2\widetilde{\mathbf{E}} \\
		\nabla\cdot\widetilde{\mathbf{E}} &= \frac{\partial}{\partial x}\left(E_0e^{-j\beta z}\right) = 0
	}
	\end{solution}

	\part The field $\widetilde{\mathbf{E}} = \hat{\mathbf{z}}\ E_0e^{-j\beta z}$
	also satisfies the Helmholtz equation. Explain why it cannot exist in a
	source-free region.
	\begin{solution}[1in]
	Its divergence is
	$\partial\widetilde{E}_z/\partial z = -j\beta E_0e^{-j\beta z} \neq 0$, but
	Gauss's law with $\widetilde{\rho}_v = 0$ requires $\nabla\cdot\widetilde{\mathbf{E}} = 0$.
	A field that varies only along $z$ cannot point along $z$ in a source-free region.
	\end{solution}

\end{parts}

\newpage
%%%%%%%%%%%%%%%%%%%%%%%%%%%%%%%%%%%%%%%%%
\question \textbf{The speed of the wave and the line.}
\begin{parts}

	\begin{minipage}[t]{0.58\linewidth}
		\part Use $\mu_0 = 4\pi\times10^{-7}$ H/m and $\epsilon_0 = 8.854\times10^{-12}$ F/m
		to find the speed of a wave in free space.
		\begin{solution}[0.9in]
		\eq{
			v_p &= \frac{1}{\sqrt{\mu_0\epsilon_0}} = \frac{1}{\sqrt{4\pi\times10^{-7}\times8.854\times10^{-12}}} \\
			&= 2.998\times10^{8}\ \text{m/s}
		}
		\end{solution}
	\end{minipage}\hfill%
	\begin{minipage}[t]{0.37\linewidth}
		\raggedleft \vspace{12pt}
		$v_p$ = \ansbox[1.3in]{$2.998\times10^{8}$ m/s}
	\end{minipage}

	\begin{minipage}[t]{0.58\linewidth}
		\part A coaxial line has $L^{\prime} = 250$ nH/m and $C^{\prime} = 100$ pF/m,
		and its dielectric has $\mu_r = 1$. Using $L^{\prime}C^{\prime} = \mu\epsilon$ for
		the coax (Lesson 30), find $\epsilon_r$ and the phase velocity.
		\begin{solution}[1.2in]
		\eq{
			\epsilon_r &= \frac{L^{\prime}C^{\prime}}{\mu_0\epsilon_0} = \frac{250\times10^{-9}\times100\times10^{-12}}{1.113\times10^{-17}} = 2.25 \\
			v_p &= \frac{1}{\sqrt{L^{\prime}C^{\prime}}} = 2.00\times10^{8}\ \text{m/s}
		}
		The dielectric is polyethylene.
		\end{solution}
	\end{minipage}\hfill%
	\begin{minipage}[t]{0.37\linewidth}
		\raggedleft \vspace{12pt}
		$\epsilon_r$ = \ansbox[1.3in]{$2.25$}

		\vspace{4pt}
		$v_p$ = \ansbox[1.3in]{$2.00\times10^{8}$ m/s}
	\end{minipage}

\end{parts}

\vspace{1.5em}
\noindent\textbf{Documentation:} \rule[-2pt]{0.6\linewidth}{0.4pt}

\end{questions}
